% Appendix E (track hw). Former Appendix B (review of representation theory, coauthor text, lightly
% edited) followed by the former main-text Sec. 3.2.6 (sub:irreps), rewritten, with the multiplicity
% computation that the paper lacked. Changes of content: jpw-review/hw/content-changes.md, part H.
\section{Representation Theory of the Twirled Channel}\label{sec:rep-theory}

\autoref{sub:fit_model} states that the Clifford twirl of the noise has two eigenvalues for a physical qutrit and up to $2542$ for the encoded qutrit. This appendix derives these numbers. It first reviews the results on representations of finite groups that the derivation uses, following~\cite{sagan2001symmetric}, and then applies them to the logical Clifford group of the $[[5,1,2]]_3$ code.

\subsection{Group Modules and Irreducible Representations}\label{app:rep-modules}

Let $G$ denote a finite group with identity element $e$.

A $G$-module is a finite dimensional vector space $V$ equipped with a group action $\cdot : G \times V \to V$ satisfying 
\begin{enumerate}
    \item \textbf{(Identity)} For all vectors $\ket{v} \in V$, the equality $e \cdot \ket{v} = \ket{v}$ is satisfied.
    \item \textbf{(Linearity)} For all linear combinations $\alpha \ket{v} + \beta \ket{w}$ and all group elements $g \in G$, the equality $g \cdot (\alpha \ket{v}+\beta\ket{w}) = \alpha (g\cdot \ket{v}) + \beta (g \cdot \ket{w})$ is satisfied.
    \item \textbf{(Compatibility)} For all group elements $g,h \in G$ and vectors $\ket{v} \in V$, the equality $(gh) \cdot \ket{v} = g\cdot (h \cdot \ket{v})$ is satisfied.
\end{enumerate}
The linearity axiom implies the existence of a map $\pi: G \to \mathrm{End}(V)$ mapping the group element $g$ to the linear endomorphism $\ket{v} \mapsto g\cdot \ket{v}$. 
The identity and compatibility axioms imply that $\pi$ is actually a group homomorphism. 

\begin{lemma}
    We call a group homomorphism $\pi : G \to \mathrm{End}(V)$ a representation of $G$. If $\pi: G \to \mathrm{End}(V)$ is a representation, then $V$ is a $G$-module under 
    \begin{align}
        g \cdot \ket{v} = \pi(g)\ket{v}.
    \end{align}
    Moreover, if $V$ is a $G$-module, then the map $\pi(g)(\ket{v}) = g\cdot \ket{v}$ is a representation of $G$.
\end{lemma}
We therefore use the terms representation of $G$ and $G$-module interchangeably.
As with all linear maps, representations are typically described by how each element of $G$ acts on each basis vector of $V$.
% Alternatively, one can describe a representation by first describing how the group $G$ acts on some set $X$ and then extending the action to the free vector space $\mathbb{C}[X]$ by linearity.

A representation can sometimes be reduced to a direct sum of smaller representations,
\begin{align}
    \pi(g) = \pi_0(g) \oplus \pi_1(g) \oplus \dots \oplus \pi_r(g) ,
\end{align}
with each $\pi_i$ a representation.
In this case, we write $\pi = \pi_0 \oplus \dots \oplus \pi_r$.
This corresponds to a decomposition of $G$-modules $V = V_0 \oplus \dots \oplus V_r$ where each $V_i$ is a $G$-module.
If no such decomposition exists, then we call the representation (or $G$-module) irreducible.
We abbreviate irreducible representation as irrep.

\begin{theorem}[Maschke's Theorem \cite{sagan2001symmetric} Theorem 1.5.3] \label{thm:maschkes-theorem}
    Every representation $\pi : G \to \mathrm{End}(V)$ decomposes into a direct sum
    \begin{align}
        \pi \cong m_0 \pi_0 \oplus m_1 \pi_1 \oplus \dots \oplus m_r \pi_r
    \end{align}
    where the $\pi_i$ are distinct irreducible representations of $G$ and $m_i \pi_i$ denotes the direct sum of $m_i$ copies of $\pi_i$. The integer $m_i$ is the multiplicity of $\pi_i$ in $\pi$.

    Similarly, if $V$ is a finite dimensional $G$-module, then 
    \begin{align}
        V = m_0 V_0 \oplus \dots \oplus m_r V_r
    \end{align}
    where the $V_i$ are distinct irreducible $G$-modules.
\end{theorem}
Maschke's theorem requires that the characteristic of the underlying field does not divide $|G|$, which holds over $\mathbb{C}$.
Its proof averages over the group, so an analog holds for compact Lie groups.
The form of Schur's lemma used below requires in addition an algebraically closed field such as $\mathbb{C}$.

\subsubsection{Schur's Lemma}

Suppose that $V$ is a $G$-module. 
We let $\mathrm{End}_G(V)$ denote all of the elements $T \in \mathrm{End}(V)$ such that for all $g \in G$
\begin{align}
    T \pi(g) = \pi(g)T.
\end{align}
Such elements $T$ are said to be invariant under the action of $G$, or $G$-linear.

\begin{lemma}[Schur's Lemma \cite{sagan2001symmetric} Theorem 1.6.5]\label{lem:schurs-lemma}
    Let $V$ and $W$ be $G$-modules and let $\pi_V$ and $\pi_W$ be the associated representations.
    We define $\mathrm{Hom}_G(V, W)$ by $T \in \mathrm{Hom}_G(V, W)$ if for all $g \in G$
    \begin{align}
        T \pi_V(g) = \pi_W(g) T.
    \end{align}
    If $V$ and $W$ are irreducible, then exactly one of the following holds.
    \begin{enumerate}
        \item The $G$-modules $V$ and $W$ are not isomorphic and $\mathrm{Hom}_G(V,W) = 0$.
        \item The $G$-modules $V$ and $W$ are isomorphic and $\mathrm{Hom}_G(V,W) \cong \mathbb{C}$ in the sense that $\mathrm{End}_G(V) = \{ \lambda I_V \mid \lambda \in \mathbb{C} \}$.
    \end{enumerate}
\end{lemma}

\begin{cor}[\cite{sagan2001symmetric} Theorem 1.7.9]\label{cor:schur-decomposition}
    Let $V$ be a $G$-module such that
    \begin{align}
        V = m_0 V_0 \oplus \dots \oplus m_r V_r
    \end{align}
    with the $V_i$ distinct irreducible modules.
    Let $d_i$ be the dimension of $V_i$.
    Then the following hold.
    \begin{enumerate}
        \item The dimension of $V$ is given by $\mathrm{dim}(V) = m_0 d_0 + \dots + m_r d_r$.
        \item \label{item:artin-wedderburn} The endomorphism algebra of $V$ is isomorphic to the direct sum
        \begin{align}
            \mathrm{End}_G(V) \cong \bigoplus_{i = 0}^r \mathrm{End}(\mathbb{C}^{m_i}).
        \end{align}
        \item The endomorphism algebra $\mathrm{End}_G(V)$ has dimension $m_0^2 + m_1^2 + \dots + m_r^2$.
    \end{enumerate}
\end{cor}

\begin{proof}
    We prove only item~\ref{item:artin-wedderburn}, the one used below.
    From the decomposition of $V$, we get the equality
    \begin{align}
        \mathrm{End}_G(V) \cong \bigoplus_{i,j = 0}^r \mathrm{Hom}_G(m_iV_i,m_jV_j).
    \end{align}
    The elements of $\mathrm{Hom}_G(m_i V_i, m_j V_j)$ can be thought of as $m_i \times m_j$ matrices whose entries are in $\mathrm{Hom}_G(V_i, V_j)$. 
    However, if $i \neq j$, then by Schur's lemma we know that $\mathrm{Hom}_G(V_i, V_j)$ only contains the zero map. 
    Thus, all of the terms in the direct sum for $i \neq j$ are zero and can be ignored. 
    When $i = j$, the elements of $\mathrm{Hom}_G(V_i,V_i) = \mathrm{End}_G(V_i)$ are all multiples of the identity map. 
    So, the elements $g_i \in \mathrm{End}_G(m_iV_i)$ can be written in the form
    \begin{align}
        g_i = M_i \otimes I_{V_i},
    \end{align}
    where $M_i \in \mathrm{End}(\mathbb{C}^{m_i})$ is an $m_i \times m_i$ matrix with complex entries. 
    % Moreover, any such matrix $M_i$ defines an element $g_i = M_i \otimes I_{V_i} \in \mathrm{End}_G(m_i V_i)$.
    A routine argument shows $\mathrm{End}_G(m_iV_i)$ is isomorphic to the matrix algebra $\mathrm{End}(\mathbb{C}^{m_i})$ under the map $g_i \mapsto M_i$.
    Finally, we conclude that $\mathrm{End}_G(V) \cong \bigoplus_{i = 0}^r \mathrm{End}_G(m_i V_i)$ and hence
    \begin{align}
        \mathrm{End}_G(V) \cong \bigoplus_{i=0}^r \mathrm{End}(\mathbb{C}^{m_i}).
    \end{align}
\end{proof}

The property of $\mathrm{End}_G(V)$ that we use is the following: if $T \in \mathrm{End}_G(V)$ and $\ket{v} \in m_i V_i$ is supported entirely in copies of the module $V_i$, then $T\ket{v} \in m_i V_i$ is also supported entirely in copies of $V_i$.
The next lemma guarantees that a twirled map is of this kind.

\begin{lemma}[Twirling] \label{lem:twirling}
    Suppose $V$ is a finite-dimensional $G$-module with associated representation $\pi$, for $G$ finite. Given any $T \in \mathrm{End}(V)$, which need not be $G$-linear, the twirled map
    \begin{align}
        T_G = \frac{1}{|G|} \sum_{g \in G} \pi(g)^{-1} T \pi(g)
    \end{align}
    is an element of $\mathrm{End}_G(V)$.
\end{lemma}

\begin{proof}
    Fix $g \in G$. Conjugating $T_G$ by $\pi(g)$ yields 
    \begin{align}
        \pi(g)^{-1} T_G \pi(g) &= \frac{1}{|G|} \sum_{h \in G} \pi(g)^{-1} \pi(h)^{-1} T \pi(h) \pi(g) \nonumber \\
        &= \frac{1}{|G|} \sum_{h \in G} \pi(hg)^{-1} T \pi(hg).
    \end{align}
    Since the right coset $Gg$ is equal to $G$, we can reindex the sum to obtain 
    \begin{align}
        \pi(g)^{-1} T_G \pi(g) &= \frac{1}{|G|} \sum_{h \in G} \pi(hg)^{-1} T \pi(hg) \nonumber \\
        &= \frac{1}{|G|} \sum_{h \in G} \pi(h)^{-1} T \pi(h) = T_G.
    \end{align}
    We conclude that $\pi(g)^{-1} T_G \pi(g) = T_G$ and hence $T_G \pi(g) = \pi(g) T_G$.
\end{proof}

\subsection{Character Theory}\label{app:rep-characters}

By item~\ref{item:artin-wedderburn} of Corollary~\ref{cor:schur-decomposition} and Lemma~\ref{lem:twirling}, the structure of a twirled map is fixed by the multiplicity of each irrep in the representation.
Character theory computes these multiplicities once the irreps are known.

\subsubsection{Characters and Conjugacy Classes}

If $\pi : G \to \mathrm{End}(V)$ is a representation, then the character $\chi_\pi : G \to \mathbb{C}$ of $\pi$ is defined by 
\begin{align}
    \chi_\pi(g) = \Tr \pi(g).
\end{align}
By the cyclic property of the trace, characters satisfy the property $\chi_\pi(h g h^{-1}) = \chi_\pi(g)$.
Characters are therefore constant on the conjugacy classes.
Such a function is called a class function.
The space of class functions has dimension equal to the number of conjugacy classes of $G$, and it carries the inner product

\begin{align}
    \langle \chi, \tau \rangle = \frac{1}{|G|} \sum_{g \in G} \chi(g)^* \tau(g).
\end{align}

\begin{lemma}[Character Orthogonality \cite{sagan2001symmetric} Theorem 1.9.3] \label{lem:character-orthogonality}
    If $\pi_1$ and $\pi_2$ are irreducible representations of the finite group $G$, then 
    \begin{align}
        \langle \chi_{\pi_1}, \chi_{\pi_2}\rangle = \begin{cases}
            1 & \pi_1 \cong \pi_2, \\
            0 & \text{otherwise}.
        \end{cases}
    \end{align}
\end{lemma}

\begin{cor}[\cite{sagan2001symmetric} Corollary 1.9.4] \label{cor:character-relations}
    Let $\pi : G \to \mathrm{End}(V)$ be a representation and $\pi = m_0 \pi_0 \oplus \dots \oplus m_r \pi_r$ be the decomposition of the representation $\pi$ into irreps. The following hold.
    \begin{enumerate}
        \item The character $\chi_{\pi}$ can be written as $\chi_{\pi} = m_0 \chi_{\pi_0} + m_1 \chi_{\pi_1} + \dots + m_r \chi_{\pi_r}$.
        \item \label{item:irrep-mult-test} The inner product $\langle \chi_{\pi}, \chi_{\pi_i}\rangle = m_i$ for all $0 \leq i \leq r$.
        \item The inner product $\langle \chi_{\pi}, \chi_{\pi} \rangle = m_0^2 + \dots + m_r^2$. In particular this is always an integer.
        \item The representation $\pi$ is irreducible if and only if $\langle \chi_{\pi}, \chi_{\pi} \rangle = 1$.
        \item If $\tau$ is another representation of $G$, then $\pi$ and $\tau$ are isomorphic if and only if $\chi_{\pi}(g) = \chi_{\tau}(g)$ for all $g \in G$.
    \end{enumerate}
\end{cor}

Lemma~\ref{lem:character-orthogonality} and Corollary~\ref{cor:character-relations} have two consequences that we use.
First, the irrep characters are orthonormal and thus linearly independent.
Together with the completeness of the irreducible characters, which span the space of class functions~\cite{sagan2001symmetric}, this shows that the irrep characters form a basis for that space and, in particular, that there are exactly as many irreps as there are conjugacy classes of $G$.
Second, item~\ref{item:irrep-mult-test} of Corollary~\ref{cor:character-relations} gives the multiplicity of an irrep in any representation of $G$ as an inner product of characters.

\subsection{The Logical Clifford Group of the \texorpdfstring{$[[5,1,2]]_3$}{[[5,1,2]]3} Code}
\label{sub:irreps}

% Former main-text Sec. 3.2.6 ("method - LRB-D.tex", lines 245-303 of the 438f6f16 mirror), rewritten. New here: the characters and the two inner products that give the multiplicities, and the table of all ten.
% All numbers of this subsection: paper-rework/scratch/lead/irreps_check.py (output in jpw-review/hw/checks/irreps_check.out): conjugation representation of the 216-element group generated by the physical Hbar and Sbar, built with sdim's gate matrices; Burnside character table of the logical action; multiplicities 301, 260, 260 (dimension 1), 560, 560, 520 (dimension 2), 820 (dimension 3), 2160, 2160, 2241 (dimension 8, the last being V_pauli); sum of m*d = 59049; the group with Xbar and Zbar has 1944 elements, m_triv = 40, m_pauli = 261; its nine-element kernel on the code space is generated by G_2 and G_4.
We now apply these results to the encoded qutrit. Let $\mathcal{H}$ be the $3^5$-dimensional Hilbert space of the five data qutrits and $\mathcal{S}$ the $3^{10}$-dimensional space of linear operators on $\mathcal{H}$, on which density operators live and noise channels act. Let $\bar{\mathcal{C}}_3$ be the group generated by the five-qutrit unitaries $\bar H$ and $\bar S$ of Eqs.~\eqref{eq:logical_H} and~\eqref{eq:logical_HS}, taken up to global phases. It has $216$ elements, and its action on the code space maps it one-to-one onto the single-qutrit Clifford group $\mathcal{C}_3$. Conjugation,
\begin{equation}
    \bar{\pi}(C)(\rho) = \bar{C} \rho\, \bar{C}^\dagger ,
\end{equation}
where $\bar C \in \bar{\mathcal{C}}_3$ is the unitary that implements $C \in \mathcal{C}_3$ on the logical qutrit, is therefore a representation of $\mathcal{C}_3$ on $\mathcal{S}$. The module of this representation is the operator space $\mathcal{S}$ and not $\mathcal{H}$.

% TODO(authors): cite the paper that lists the conjugacy classes of the single-qutrit Clifford group (was a red author note in the text; jpw-review/OPEN-QUESTIONS.md, item 7).
Up to global phases, $\mathcal{C}_3$ has $10$ conjugacy classes, of sizes $1$, $8$, $9$, $12$, $12$, $24$, $24$, $36$, $36$ and $54$, and hence $10$ irreps (\autoref{tab:irreps}). Two of them matter here. The first is the trivial irrep $V_{\mathrm{triv}}$. The second is the eight-dimensional irrep $V_{\mathrm{pauli}}$ in which $\mathcal{C}_3$ conjugates the eight traceless Pauli operators of one qutrit. The operators of one qutrit decompose as $V_{\mathrm{triv}} \oplus V_{\mathrm{pauli}}$, which is the origin of the two eigenvalues of Eq.~\eqref{eq:twirl_irrep}. By Theorem~\ref{thm:maschkes-theorem}, the operators of five qutrits decompose as
\begin{equation}
    \mathcal{S} \cong m_{\mathrm{triv}} V_{\mathrm{triv}} \oplus m_{\mathrm{pauli}} V_{\mathrm{pauli}} \oplus \bigoplus_{i = 2}^{9} m_i V_i ,
    \label{eq:irrep-decomposition}
\end{equation}
where $V_2, \dots, V_9$ are the other eight irreps.

The multiplicities follow from item~\ref{item:irrep-mult-test} of Corollary~\ref{cor:character-relations}. Conjugation by a unitary $U$ acts on the operator space as $U \otimes U^{*}$, whose trace is $\lvert \Tr U \rvert^2$, so the character of $\bar\pi$ is $\chi_{\bar\pi}(C) = \lvert \Tr \bar C \rvert^2$. In the same way the character of one-qutrit conjugation is $\lvert \Tr C \rvert^2$, with $C$ taken as a $3 \times 3$ unitary, and the character of $V_{\mathrm{pauli}}$ is $\chi_{\mathrm{pauli}}(C) = \lvert \Tr C \rvert^2 - 1$. The two multiplicities are therefore
\begin{equation}
  \begin{aligned}
    m_{\mathrm{triv}} &= \frac{1}{216} \sum_{C \in \mathcal{C}_3} \lvert \Tr \bar C \rvert^2 = 301 , \\
    m_{\mathrm{pauli}} &= \frac{1}{216} \sum_{C \in \mathcal{C}_3} \left( \lvert \Tr C \rvert^2 - 1 \right) \lvert \Tr \bar C \rvert^2 = 2241 ,
  \end{aligned}
  \label{eq:multiplicities}
\end{equation}
which we evaluated numerically from the $216$ five-qutrit unitaries. \autoref{tab:irreps} lists the multiplicities of all ten irreps, for which we computed the remaining characters from the multiplication table of the group. They satisfy $\sum_i m_i \dim V_i = 3^{10}$. For one physical qutrit, $m_{\mathrm{triv}} = m_{\mathrm{pauli}} = 1$ and all other multiplicities vanish.

\begin{table*}[tp]
  \centering
  \begin{tabular}{l|cccccccccc}
    \hline
    Irrep & $V_{\mathrm{triv}}$ & $V_2$ & $V_3$ & $V_4$ & $V_5$ & $V_6$ & $V_7$ & $V_8$ & $V_9$ & $V_{\mathrm{pauli}}$ \\
    \hline
    Dimension & 1 & 1 & 1 & 2 & 2 & 2 & 3 & 8 & 8 & 8 \\
    Multiplicity & 301 & 260 & 260 & 560 & 560 & 520 & 820 & 2160 & 2160 & 2241 \\
    \hline
  \end{tabular}
  \caption{The ten irreps of the single-qutrit Clifford group $\mathcal{C}_3$ and their multiplicities in the conjugation representation $\bar\pi$ of the group generated by $\bar H$ and $\bar S$ on the $3^{10}$-dimensional operator space of the five data qutrits.}
  \label{tab:irreps}
\end{table*}

The simulated gate set generates a larger group. It applies the logical Pauli operators as $\bar X = X_0 X_3$ and $\bar Z = Z_0 Z_2$ (Eq.~\eqref{eq:logicalPauli}), which differ from the elements of $\bar{\mathcal{C}}_3$ with the same logical action by the stabilizers $G_4$ and $G_2$. The group generated by $\bar H$, $\bar S$, $\bar X$ and $\bar Z$ thus consists of the products of the elements of $\bar{\mathcal{C}}_3$ with the nine products of powers of $G_2$ and $G_4$. It has $1944$ elements and acts on the code space as $\bar{\mathcal{C}}_3$ does. The same computation for this group gives $m_{\mathrm{triv}} = 40$ and $m_{\mathrm{pauli}} = 261$.

These multiplicities bound the number of eigenvalues that a decay can contain. Let $\mathrm{End}_{\mathcal{C}_3}(\mathcal{S})$ be the set of linear maps on $\mathcal{S}$ that commute with $\bar\pi$. By item~\ref{item:artin-wedderburn} of Corollary~\ref{cor:schur-decomposition} it is isomorphic, as an algebra, to a direct sum of matrix algebras,
\begin{align}
    \mathrm{End}_{\mathcal{C}_3}(\mathcal{S}) &\cong \mathrm{Mat}_{m_{\mathrm{triv}}}(\mathbb{C}) \otimes I_{V_{\mathrm{triv}}} \nonumber \\
    &\quad \oplus \mathrm{Mat}_{m_{\mathrm{pauli}}}(\mathbb{C})\otimes I_{V_{\mathrm{pauli}}} \nonumber \\
    &\quad \oplus \bigoplus_{i = 2}^9 \mathrm{Mat}_{m_i} (\mathbb{C}) \otimes I_{V_i} ,
\end{align}
so that such a map cannot take a vector of $m_i V_i$ to a vector with a component in $m_j V_j$ for $j \neq i$. The twirled error channel $\Lambda_{\mathrm{eff}}$ is a map of this kind by Lemma~\ref{lem:twirling}, and in a suitable basis it reads
\begin{equation}
    \Lambda_{\mathrm{eff}} = M_{\mathrm{triv}}\otimes I_{V_{\mathrm{triv}}} \oplus M_{\mathrm{pauli}} \otimes I_{V_{\mathrm{pauli}}} \oplus M ,
\end{equation}
where $M_{\mathrm{triv}}$ and $M_{\mathrm{pauli}}$ are square matrices of sizes $m_{\mathrm{triv}}$ and $m_{\mathrm{pauli}}$ and $M$ collects the other eight irreps. If $M_{\mathrm{triv}}$ and $M_{\mathrm{pauli}}$ are diagonalizable, a further change of basis gives
\begin{equation}
  \begin{aligned}
    \Lambda_{\mathrm{eff}} &= \Lambda_{\mathrm{triv}} \oplus \Lambda_{\mathrm{pauli}} \otimes I_{V_{\mathrm{pauli}}} \oplus M , \\
    \Lambda_{\mathrm{eff}}^{\,t} &= \Lambda_{\mathrm{triv}}^{\,t} \oplus \Lambda_{\mathrm{pauli}}^{\,t} \otimes I_{V_{\mathrm{pauli}}} \oplus M^{t} ,
  \end{aligned}
  \label{eq:twirl-logical}
\end{equation}
with $\Lambda_{\mathrm{triv}}$ and $\Lambda_{\mathrm{pauli}}$ the diagonal matrices of their eigenvalues and $t$ the number of applications of the channel.

The block $M$ does not act on a state of the code space. The operators supported on the code space transform under $\bar\pi$ as the operators of one qutrit, so a code state such as $\ket{\bar{+}}\bra{\bar{+}}$ decomposes as $\rho = \rho_{\mathrm{triv}} \oplus \rho_{\mathrm{pauli}} \oplus 0$ with $\rho_{\mathrm{triv}} \in m_{\mathrm{triv}} V_{\mathrm{triv}}$ and $\rho_{\mathrm{pauli}} \in m_{\mathrm{pauli}}V_{\mathrm{pauli}}$. By Eq.~\eqref{eq:twirl-logical} its evolution under repeated application of $\Lambda_{\mathrm{eff}}$ contains only the powers of the diagonal elements of $\Lambda_{\mathrm{triv}}$ and $\Lambda_{\mathrm{pauli}}$. There are $m_{\mathrm{triv}} + m_{\mathrm{pauli}} = 301 + 2241 = 2542$ of them for $\bar{\mathcal{C}}_3$, and $40 + 261 = 301$ for the group that the simulated gate set generates, against $1 + 1 = 2$ for a physical qutrit. If $M_{\mathrm{triv}}$ or $M_{\mathrm{pauli}}$ is not diagonalizable, its Jordan form takes the place of the diagonal matrix, and a Jordan block of size $s$ with eigenvalue $\lambda$ contributes the terms $\binom{t}{j}\lambda^{t-j}$ with $j < s$ in place of the single power $\lambda^t$.
